\chapter{How Conjectures Are Born and Die}\label{conjectures:ch:born}

\begin{center}\itshape
There aren't enough small numbers to meet the many demands made of them.
\upshape(Richard K.~Guy, \emph{The strong law of small numbers}, 1988)
\end{center}

In 1919 George P\'olya noticed something about the numbers below a few thousand. Count the integers $n\le x$ that have an even number of prime factors
(counted with multiplicity, so $12=2\cdot2\cdot3$ has three), and those that have an odd number. P\'olya checked that the odd ones were never outnumbered: for every
$x\ge2$ he tried, at least half of the numbers up to $x$ had an odd number of prime factors. Others pushed the check to a million, and it held. In 1958 Brian
Haselgrove proved that it fails somewhere, without saying where; in 1960 Sherman Lehman found a failure at $906\,180\,359$; and in 1980 Minoru Tanaka found the
first one, at $906\,150\,257$. For more than nine hundred million numbers the pattern is perfect, and then it breaks.

This book is about open problems: the twin primes, Goldbach, the Riemann hypothesis, Collatz, and a dozen others. Before we meet them, this chapter asks a
question that hangs over all of them. When a pattern has been checked for every number up to a billion, or up to $10^{18}$, why should we believe it? We
look at conjectures that died, some of them after a very long life, and at what killed them. We prove two theorems that explain the deaths: a theorem of
Landau about positive Dirichlet series, which shows that P\'olya's pattern, had it been true, would have proved the Riemann hypothesis; and a probabilistic
explanation of the Borwein integrals, a sequence of seven integrals equal to $\frac\pi2$ whose eighth member misses $\frac\pi2$ by $2.3\times10^{-11}$. The lesson
is not that conjectures are worthless. It is that a conjecture is only as good as the reason behind it, and the rest of this book is about those reasons.

\begin{think}
The polynomial $n^2+n+41$ is prime for $n=0,1,2,\dots,39$. How much would you bet that it is prime for $n=40$? Before computing, look at $n=41$.
\end{think}

\section{A small gallery of broken patterns}

Some patterns die young. Euler noticed that $n^2+n+41$ is prime for $n=0,\dots,39$; at $n=40$ it is $40^2+40+41=41^2$, and at $n=41$ it is obviously divisible
by $41$. No non-constant polynomial with integer coefficients can be prime for every $n$ (Exercise~\ref{conjectures:ex:euler41}), so this one had to fail somewhere; the
surprise is only how late.

Fermat believed that all the numbers $F_n=2^{2^n}+1$ are prime. They are for $n=0,1,2,3,4$: $3,5,17,257,65537$. In 1732 Euler found that
$F_5=4\,294\,967\,297=641\times6\,700\,417$ (Exercise~\ref{conjectures:ex:fermat641}), and no prime $F_n$ beyond $F_4$ has been found since.

Draw $n$ points on a circle in general position and join them all by chords. The regions inside the circle number $1,2,4,8,16$ for $n=1,\dots,5$, and the reader
who guesses $32$ for six points finds $31$ (Exercise~\ref{conjectures:ex:moser}). The numbers $31$, $331$, $3331$, $33331$, $333331$, $3333331$, $33333331$ are all prime, and
$333333331=17\times19\,607\,843$ (Exercise~\ref{conjectures:ex:threes}). Guy collected dozens of such stories in a paper whose title is this chapter's epigraph, and drew the
moral: small numbers are few, and they are asked to carry too many patterns, so coincidences among them are inevitable.

Other patterns live long. In a letter to Euler of 18 November 1752, Goldbach suggested that every odd number is a prime plus twice a square, $p+2a^2$ (with $1$
counted as a prime and $a=0$ allowed, as was usual then). A century later, in 1856, Moritz Stern and his students checked the odd numbers up to $9000$ and
found two exceptions, $5777=53\cdot109$ and $5993$; none is known beyond them. Euler himself guessed that at least $k$ $k$-th powers are needed to make a
$k$-th power (for $k=3$ this is Fermat's last theorem for cubes). In 1966 Leon Lander and Thomas Parkin, searching on a CDC~6600 computer, found
\[
27^5+84^5+110^5+133^5=144^5,
\]
and published it in a paper of two sentences. For fourth powers the counterexamples are much rarer: Noam Elkies found the first in 1988,
\[
2\,682\,440^4+15\,365\,639^4+18\,796\,760^4=20\,615\,673^4,
\]
from an elliptic curve that produces infinitely many, and Roger Frye then found the smallest, $95\,800^4+217\,519^4+414\,560^4=422\,481^4$, by a computer search
guided by Elkies' method. A reader with Python can check all three identities in a second; nobody found them by hand in two centuries.

\section{P\'olya's problem}

The \emph{Liouville function} is $\lambda(n)=(-1)^{\Omega(n)}$, where $\Omega(n)$ is the number of prime factors of $n$ counted with multiplicity, so
$\lambda(1)=1$, $\lambda(p)=-1$, $\lambda(12)=-1$. It is completely multiplicative: $\lambda(mn)=\lambda(m)\lambda(n)$ for all $m,n$. Its summatory function
\[
L(x)=\sum_{n\le x}\lambda(n)
\]
is the number of $n\le x$ with an even number of prime factors minus the number with an odd number. P\'olya's observation was that $L(x)\le0$ for all $x\ge2$, and
he asked whether it is always so. The left half of Figure~\ref{conjectures:fig:liouville} shows $L(x)/\sqrt x$ up to $10^9$: it hovers around $-\frac12$ and comes near $0$
only rarely. Haselgrove's 1958 proof that $L(x)>0$ for some $x$ was indirect, through the zeros of $\zeta$; Lehman's 1960 computation found an explicit example,
and Tanaka's 1980 search found the first: $L(906\,150\,257)=1$. The positive stretch, shown on the right of Figure~\ref{conjectures:fig:liouville}, lasts with a few dips
until $906\,488\,079$, and its highest point is $L(906\,316\,571)=829$. We recomputed all these numbers with a segmented sieve, in about a minute (a numerical calculation
of this chapter does the same in Python).

\begin{figure}[ht]
\centering
\fig{conjectures-ch01-liouville}
\caption{Left: $L(x)/\sqrt x$ for $10\le x\le10^9$, on a logarithmic scale in $x$. It stays below $0$, around $-\frac12$, except on a short stretch that the
picture is too coarse to show. Right: $L(x)$ for $905.9\times10^6\le x\le906.7\times10^6$. It becomes positive for the first time at $x=906\,150\,257$ and
stays positive, apart from a few dips, until $906\,488\,079$.}
\label{conjectures:fig:liouville}
\end{figure}

Why did P\'olya care? The generating Dirichlet series of $\lambda$ is a quotient of two zeta functions. For $\operatorname{Re}s>1$ the Euler product
(Volume~5, Theorem~12.3) gives
\[
\sum_{n\ge1}\frac{\lambda(n)}{n^s}=\prod_p\Bigl(1-\frac1{p^s}+\frac1{p^{2s}}-\cdots\Bigr)=\prod_p\frac1{1+p^{-s}}=\prod_p\frac{1-p^{-s}}{1-p^{-2s}}=\frac{\zeta(2s)}{\zeta(s)}.
\]
The zeros of $\zeta$ in the critical strip are poles of the right side, so the size of $L(x)$ is governed by them, as the size of $\psi(x)-x$ is in Volume~6,
Chapter~7. P\'olya saw that a sign condition alone, with no size condition at all, would force the Riemann hypothesis. The tool is a theorem of Landau.

\section{Landau's theorem}

\begin{theorem}[Landau, 1905]\label{conjectures:thm:landau}
Let $A:[1,\infty[\to\R$ be bounded on every bounded interval, piecewise continuous, and such that $A(x)\ge0$ for all $x\ge x_0$. Suppose that the integral
\[
F(s)=\int_1^\infty A(x)\,x^{-s-1}\dd x
\]
converges for some real $s$, and let $\sigma_c$ be the infimum of the real $s$ for which it converges. Then $F$ is holomorphic on $\operatorname{Re}s>\sigma_c$,
and there is no holomorphic function on a disc around $\sigma_c$ that agrees with $F$ on the part of the disc where $\operatorname{Re}s>\sigma_c$.
\end{theorem}

In words: a Mellin transform of an eventually positive function always has a singularity at the real point where its convergence stops. It cannot stop
converging for no visible reason.

\begin{proof}
The integral over $[1,x_0]$ is an entire function of $s$ (differentiate under the integral sign over a compact interval), so it changes neither the
convergence nor the singularities, and we may assume $A\ge0$ everywhere. For real $\sigma$ and complex $s$ with $\operatorname{Re}s\ge\sigma$ we have
$|A(x)x^{-s-1}|=A(x)x^{-\operatorname{Re}s-1}\le A(x)x^{-\sigma-1}$, so convergence at a real $\sigma$ gives absolute convergence on the closed half-plane
$\operatorname{Re}s\ge\sigma$. Moreover, for every $k\ge0$ and every $\varepsilon>0$, $(\log x)^k\le C_{k,\varepsilon}x^\varepsilon$ on $[1,\infty[$. So for
$\operatorname{Re}s\ge\sigma+\varepsilon$ with $\sigma>\sigma_c$, the integrands $A(x)(\log x)^kx^{-s-1}$ are dominated by $C_{k,\varepsilon}A(x)x^{-\sigma-1}$,
which is integrable. Differentiation under the integral sign then shows that $F$ is holomorphic on $\operatorname{Re}s>\sigma_c$, with
\[
F^{(k)}(s)=\int_1^\infty A(x)\,(-\log x)^k\,x^{-s-1}\dd x.
\]

Suppose now that $G$ is holomorphic on a disc $D(\sigma_c,\delta)$ and equals $F$ where both are defined. Then $F$ and $G$ glue to a holomorphic function $H$ on
$U=\{\operatorname{Re}s>\sigma_c\}\cup D(\sigma_c,\delta)$. Fix the real point $a=\sigma_c+1$ and a small $\eta>0$, and let $R=1+\eta$. A point $s$ of the disc
$D(a,R)$ with $\operatorname{Re}s\le\sigma_c$ satisfies $\operatorname{Re}s>a-R=\sigma_c-\eta$ and $(\operatorname{Im}s)^2<R^2-1=2\eta+\eta^2$, so
$|s-\sigma_c|^2<\eta^2+2\eta+\eta^2$, which is less than $\delta^2$ once $\eta$ is small enough. For such an $\eta$, the disc $D(a,R)$ lies in $U$, and the Taylor
series of $H$ at $a$ converges on it. At the real point $s=\sigma_c-\frac\eta2$, inside $D(a,R)$, it reads
\[
H(s)=\sum_{k\ge0}\frac{F^{(k)}(a)}{k!}(s-a)^k=\sum_{k\ge0}\frac{(a-s)^k}{k!}\int_1^\infty A(x)(\log x)^kx^{-a-1}\dd x,
\]
since $(s-a)^k(-\log x)^k=(a-s)^k(\log x)^k$. Every term is nonnegative, because $a-s>0$, $A\ge0$ and $\log x\ge0$. So we may exchange the sum and the integral
(monotone convergence), and the exponential series sums to $x^{a-s}$:
\[
H(s)=\int_1^\infty A(x)\,x^{-a-1}\sum_{k\ge0}\frac{((a-s)\log x)^k}{k!}\dd x=\int_1^\infty A(x)\,x^{-s-1}\dd x.
\]
The left side is finite, so the integral converges at the real point $s=\sigma_c-\frac\eta2<\sigma_c$. This contradicts the definition of $\sigma_c$.
\end{proof}

Positivity is essential. The alternating series $\eta(s)=\sum(-1)^{n-1}n^{-s}$ converges exactly for real $s>0$, yet it extends to an entire function
(Exercise~\ref{conjectures:ex:etaentire}): its convergence stops at $0$ for no visible reason, because cancellation, not size, made it converge.

\section{P\'olya's pattern would have proved the Riemann hypothesis}

We need three facts about $\zeta$ from Volume~5. It extends to a holomorphic function on $\operatorname{Re}s>0$ except for a simple pole at $s=1$
(Volume~5, Proposition~12.6). It has no zeros on $\operatorname{Re}s>1$, by the Euler product. And its zeros in the strip $0<\operatorname{Re}s<1$ are symmetric
with respect to the line $\operatorname{Re}s=\frac12$ (Volume~6, Proposition~6.2). We also need one fact about real values.

\begin{lemma}\label{conjectures:lem:zetaneg}
For real $0<\sigma<1$, $\zeta(\sigma)<0$.
\end{lemma}
\begin{proof}
For $\operatorname{Re}s>1$, $\eta(s)=\sum(-1)^{n-1}n^{-s}=\zeta(s)-2\cdot2^{-s}\zeta(s)=(1-2^{1-s})\zeta(s)$. The series $\eta$ converges on $\operatorname{Re}s>0$
and is holomorphic there (Volume~5, Chapter~12), so by the identity theorem the equality $\eta(s)=(1-2^{1-s})\zeta(s)$ holds on $\operatorname{Re}s>0$, $s\ne1$. For real
$0<\sigma<1$, grouping the terms in pairs gives $\eta(\sigma)=(1-2^{-\sigma})+(3^{-\sigma}-4^{-\sigma})+\cdots>0$, while $1-2^{1-\sigma}<0$. So $\zeta(\sigma)<0$.
\end{proof}

\begin{theorem}[P\'olya, 1919]\label{conjectures:thm:polyarh}
If $L(x)\le0$ for all sufficiently large $x$, then the Riemann hypothesis holds: every zero of $\zeta$ with $0<\operatorname{Re}s<1$ has real part $\frac12$.
\end{theorem}
\begin{proof}
Since $|L(x)|\le x$, partial summation gives, for $\operatorname{Re}s>1$,
\[
\frac{\zeta(2s)}{\zeta(s)}=\sum_{n\ge1}\frac{\lambda(n)}{n^s}=s\int_1^\infty L(x)\,x^{-s-1}\dd x.
\]
(Write $n^{-s}=s\int_n^\infty x^{-s-1}\dd x$ and exchange the sum and the integral, which is legitimate since $\sum_n|\lambda(n)|\int_n^\infty|x^{-s-1}|\dd x<\infty$.)
Apply Theorem~\ref{conjectures:thm:landau} to $A=-L$, which is nonnegative for large $x$. The integral $F(s)=\int_1^\infty A(x)x^{-s-1}\dd x$ converges for $s>1$, so $\sigma_c\le1$,
and $F(s)=G(s)$ for $\operatorname{Re}s>1$, where
\[
G(s)=-\frac{\zeta(2s)}{s\,\zeta(s)}.
\]
The function $G$ is meromorphic on $\operatorname{Re}s>\frac12$: there $\zeta(2s)$ is holomorphic (as $\operatorname{Re}2s>1$) and never $0$, so the poles of $G$ are
exactly the zeros of $\zeta$ in $\operatorname{Re}s>\frac12$. On the real half-line $]\frac12,\infty[$, $G$ is holomorphic: $\zeta$ has no zero there, by
Lemma~\ref{conjectures:lem:zetaneg} on $]\frac12,1[$ and by the Euler product on $]1,\infty[$, and at $s=1$ the pole of $\zeta$ only makes $G$ vanish.

Suppose $\sigma_c>\frac12$. Then $G$ is holomorphic on a disc around the real point $\sigma_c$ and agrees with $F$ on its right part, by the identity theorem (both
are holomorphic on $\operatorname{Re}s>\sigma_c$ away from the poles of $G$, and agree on $\operatorname{Re}s>1$). This contradicts Theorem~\ref{conjectures:thm:landau}. So
$\sigma_c\le\frac12$, and $F$ is holomorphic on $\operatorname{Re}s>\frac12$. If $\zeta(\rho)=0$ with $\operatorname{Re}\rho>\frac12$, then $G=F$ near $\rho$ except
at $\rho$, and $G$ has a pole at $\rho$ while $F$ is bounded near $\rho$: impossible. So $\zeta$ has no zeros with $\operatorname{Re}s>\frac12$, and by the
symmetry about the critical line it has none with $0<\operatorname{Re}s<\frac12$ either.
\end{proof}

So P\'olya's pattern was not an idle curiosity: its truth would have settled the most famous problem in mathematics. Its failure says nothing against the
Riemann hypothesis, and in fact the same function explains why the pattern was so convincing. The pole of $\zeta(2s)$ at $s=\frac12$ gives $\frac{\zeta(2s)}{s\zeta(s)}$ a pole at $\frac12$ with
residue $\frac1{\zeta(1/2)}=-0.6848\ldots$ (Exercise~\ref{conjectures:ex:residue}). If the Riemann hypothesis holds and the zeros are simple, the contour argument behind the
explicit formula of Volume~6, Chapter~7 suggests
\[
L(x)\approx\frac{\sqrt x}{\zeta(\frac12)}+\sum_\rho(\text{oscillating terms of size about }\sqrt x/|\rho|).
\]
The squares, which all have $\lambda(n)=1$, pull $L(x)$ down by about $0.68\sqrt x$. The oscillating terms come from the zeros $\rho=\frac12+i\gamma$ and have
amplitudes that decrease like $\frac1{|\rho\zeta'(\rho)|}$; to beat the bias, many of them must line up in phase, and that happens for the first time near
$9\times10^8$. Figure~\ref{conjectures:fig:liouville} shows the bias and the rare excursion. The same mechanism makes primes $3\bmod4$ lead primes $1\bmod4$ most of the time
(Chebyshev's bias, Volume~6, Chapter~7), and it makes $\pi(x)<\operatorname{li}(x)$ for every $x$ ever computed.

\begin{remark}[Mertens, Littlewood and Skewes]
Two other famous patterns died the same way. In 1885 Thomas Stieltjes wrote to Charles Hermite that he could prove $M(x)=\sum_{n\le x}\mu(n)=O(\sqrt x)$, and
in 1897 Franz Mertens conjectured the stronger $|M(x)|<\sqrt x$ for $x>1$, which would imply the Riemann hypothesis (Volume~6, Exercise~7.3). In 1985 Andrew
Odlyzko and Herman te Riele disproved it: $\limsup M(x)/\sqrt x>1.06$ and $\liminf M(x)/\sqrt x<-1.009$. No $x$ with $|M(x)|>\sqrt x$ is known; Pintz (1987) proved
that one exists below an explicit, astronomically large bound, and Hurst (2018) computed $M(x)$ for all $x\le10^{16}$ without finding one. Likewise
$\pi(x)<\operatorname{li}(x)$ for every $x$ ever computed, yet Littlewood proved in 1914 that the difference changes sign infinitely often (Volume~6, Chapter~5).
Skewes showed in 1933, assuming the Riemann hypothesis, that a sign change occurs below $e^{e^{e^{79}}}$, and Bays and Hudson (2000) showed that there are sign
changes near $1.398\times10^{316}$.
\end{remark}

\begin{think}
Both $L(x)$ and $\pi(x)-\operatorname{li}(x)$ have a bias of size $\sqrt x$ (up to logarithms) and oscillations of the same size. Why is the excursion of $L$
found below $10^9$, while that of $\pi-\operatorname{li}$ is far beyond $10^{300}$? Compare the size of the bias with the size of the first oscillating term in each
case.
\end{think}

\section{The Borwein integrals}

In 2001 David and Jonathan Borwein published a series of integrals that has fooled many computer algebra systems and their users. Write
$\operatorname{sinc}x=\frac{\sin x}x$ (with $\operatorname{sinc}0=1$). Then
\[
\int_0^\infty\operatorname{sinc}x\dd x=\frac\pi2,\qquad\int_0^\infty\operatorname{sinc}x\operatorname{sinc}\frac x3\dd x=\frac\pi2,\qquad
\int_0^\infty\operatorname{sinc}x\operatorname{sinc}\frac x3\operatorname{sinc}\frac x5\dd x=\frac\pi2,
\]
and so on up to the factor $\operatorname{sinc}\frac x{13}$: seven integrals, all equal to $\frac\pi2$. The eighth,
\[
\int_0^\infty\operatorname{sinc}x\operatorname{sinc}\frac x3\cdots\operatorname{sinc}\frac x{15}\dd x
=\frac{467\,807\,924\,713\,440\,738\,696\,537\,864\,469}{935\,615\,849\,440\,640\,907\,310\,521\,750\,000}\,\pi,
\]
is smaller than $\frac\pi2$ by about $2.31\times10^{-11}$. The explanation is a piece of probability.

\begin{theorem}\label{conjectures:thm:borwein}
Let $a_1,\dots,a_n>0$, let $U_1,\dots,U_n$ be independent and uniform on $[-1,1]$, and let $S=a_1U_1+\dots+a_nU_n$. Then
\[
\int_0^\infty\operatorname{sinc}x\prod_{k=1}^n\operatorname{sinc}(a_kx)\dd x=\frac\pi2\,\mathbb P\bigl(|S|<1\bigr).
\]
In particular the integral equals $\frac\pi2$ if $a_1+\dots+a_n\le1$, and is strictly smaller if $a_1+\dots+a_n>1$.
\end{theorem}
\begin{proof}
For a uniform $U$ on $[-1,1]$ and real $t$, $\mathbb E[\cos(tU)]=\frac12\int_{-1}^1\cos(tu)\dd u=\operatorname{sinc}t$. Since $S$ is symmetric,
$\mathbb E[\sin(xS)]=0$, and by independence
\[
\mathbb E[\cos(xS)]=\operatorname{Re}\mathbb E[e^{ixS}]=\operatorname{Re}\prod_k\mathbb E[e^{ixa_kU_k}]=\prod_k\operatorname{sinc}(a_kx).
\]
For $T>0$, Fubini's theorem on $[0,T]\times\Omega$ (the integrand is bounded by $1$) gives
\[
\int_0^T\operatorname{sinc}x\prod_k\operatorname{sinc}(a_kx)\dd x=\mathbb E\int_0^T\frac{\sin x\cos(xS)}x\dd x=\frac12\,\mathbb E\bigl[D_T(1+S)+D_T(1-S)\bigr],
\]
where $D_T(c)=\int_0^T\frac{\sin(cx)}x\dd x$, by $2\sin x\cos(xS)=\sin((1+S)x)+\sin((1-S)x)$. Substituting $u=|c|x$, $D_T(c)=\operatorname{sgn}(c)\int_0^{|c|T}\frac{\sin u}u\dd u$.
The function $t\mapsto\int_0^t\frac{\sin u}u\dd u$ is continuous on $[0,\infty[$ and tends to $\frac\pi2$ (Volume~5, Theorem~6.11), so it is bounded, say by $B$.
Hence $|D_T(c)|\le B$ for all $c$ and $T$, and $D_T(c)\to\frac\pi2\operatorname{sgn}(c)$ as $T\to\infty$. By dominated convergence,
\[
\int_0^\infty\operatorname{sinc}x\prod_k\operatorname{sinc}(a_kx)\dd x=\frac\pi4\,\mathbb E\bigl[\operatorname{sgn}(1+S)+\operatorname{sgn}(1-S)\bigr].
\]
The bracket equals $2$ if $|S|<1$, $1$ if $|S|=1$ and $0$ if $|S|>1$. Since $S$ has a density when $n\ge1$, $\mathbb P(|S|=1)=0$, and the value is
$\frac\pi2\mathbb P(|S|<1)$. Finally $|S|\le a_1+\dots+a_n$, so $\mathbb P(|S|<1)=1$ when the sum is at most $1$. When the sum exceeds $1$, the event
$S>1$ contains a small box of positive volume around $(1,\dots,1)$, so $\mathbb P(|S|<1)<1$.
\end{proof}

For the Borwein integrals, $a_k=\frac1{2k+1}$. The sum $\frac13+\frac15+\dots+\frac1{13}=\frac{43\,024}{45\,045}=0.955\ldots$ is less than $1$, but adding $\frac1{15}$
gives $1.0218\ldots$ (Figure~\ref{conjectures:fig:borwein}). The pattern breaks exactly when the sum crosses $1$, and the deficit is tiny because $S$ can exceed $1$ only
when all seven $U_k$ are very close to $\pm1$ with the same sign.

\begin{corollary}\label{conjectures:cor:borweindeficit}
Let $\varepsilon=\sum_{k=1}^7\frac1{2k+1}-1=\frac{982}{45\,045}$. Then
\[
\frac\pi2-\int_0^\infty\prod_{k=0}^7\operatorname{sinc}\frac x{2k+1}\dd x=\pi\,\frac{3\cdot5\cdot7\cdots15}{7!\,2^7}\,\varepsilon^7\approx2.31\times10^{-11}.
\]
\end{corollary}
\begin{proof}
By Theorem~\ref{conjectures:thm:borwein} and symmetry, the deficit is $\frac\pi2\mathbb P(|S|\ge1)=\pi\,\mathbb P(S\ge1)$. Put $V_k=1-U_k$, independent and uniform on $[0,2]$. Then
$S=\sum a_k-\sum a_kV_k$, and $S\ge1$ means $\sum a_kV_k\le\varepsilon$. This forces $V_k\le\frac\varepsilon{a_k}\le15\varepsilon<2$ for every $k$, so the event is
the simplex $\{v_k\ge0,\ \sum a_kv_k\le\varepsilon\}$, which lies inside $[0,2]^7$. Its volume is $\frac{\varepsilon^7}{7!\,a_1\cdots a_7}$, and the density of
$(V_1,\dots,V_7)$ is $2^{-7}$. Multiplying, $\mathbb P(S\ge1)=\frac{\varepsilon^7}{7!\,2^7a_1\cdots a_7}$, which is the formula. Computed in exact arithmetic, it
agrees with the Borweins' fraction above.
\end{proof}

\begin{figure}[ht]
\centering
\fig{conjectures-ch01-borwein}
\caption{Left: the densities of $S_n=\frac{U_1}3+\dots+\frac{U_n}{2n+1}$ for $n=1,2,3,6$, computed by convolution; all live inside $[-1,1]$, so the first seven
Borwein integrals equal $\frac\pi2$. Right: the partial sums $\frac13+\dots+\frac1{2n+1}$. The seventh, $1.0218$, is the first to exceed $1$: from then on $S_n$
can leave $[-1,1]$, with probability $7.4\times10^{-12}$ on each side for $n=7$.}
\label{conjectures:fig:borwein}
\end{figure}

\section{Why we believe what we believe}

These deaths share a pattern. P\'olya's conjecture, Mertens' and the inequality $\pi(x)<\operatorname{li}(x)$ were each the visible part of a bias of size about
$\sqrt x$ fighting oscillations of the same size, and an oscillation needs many zeros of $\zeta$ in phase to win, which takes astronomically long. The Borwein
pattern was a threshold, $\sum a_k\le1$, crossed late because $\sum\frac1{2k+1}$ grows like $\frac12\log n$, slowly. Euler's sum of powers conjecture was a
heuristic count (``$k$ powers give $k$ degrees of freedom'') that ignored the structure of special varieties, and it is that structure, an elliptic curve, that
Elkies used. In each case, the evidence was large, and the reason was missing or wrong.

The conjectures in the rest of this book come with better reasons. Most of them come from a probabilistic model: Cram\'er's model of the primes as random
numbers with probability $\frac1{\log n}$ (Volume~6, Definition~10.1), corrected for divisibility by small primes, predicts the number of twin primes, prime
triples and Goldbach representations with the right constants, and these predictions have matched every computation for a century. A model that predicts
constants, not just patterns, is much harder to fool. But Volume~6, Chapter~10 also shows Maier's theorem, where even the corrected model fails. A good
conjecture is a precise prediction that has survived attempts to kill it, by computation and by theory; it is never a proof.

\section{Reading a paper: Guy's strong law}

Richard Guy's \emph{The strong law of small numbers} (American Mathematical Monthly, 1988) is a gallery of 35 examples, presented as a quiz: for each pattern,
the reader is asked whether it continues. Some do, some do not, and several were open when Guy wrote. It needs no prerequisites, and it is the right first paper
for this book.
\begin{itemize}
\item \emph{What it says.} Two ``laws'': there are not enough small numbers for the demands made of them, and (in its 1990 sequel in Mathematics Magazine) when
two numbers look equal, it ain't necessarily so.
\item \emph{How to read it.} Make your own guess for each example before reading the answer, and keep a score. Then check the examples you got wrong with
Python.
\item \emph{The idea worth stealing.} When a sequence matches a familiar one for its first few terms, ask how many familiar sequences have those first terms:
the answer is usually many, and the On-Line Encyclopedia of Integer Sequences will list them.
\end{itemize}

\section{Open doors}

\begin{itemize}
\item Is $5777$, $5993$ the complete list of exceptions to Goldbach's ``other conjecture''? Every computation says yes. Estimate, with Cram\'er's model, the expected
number of representations $n=p+2a^2$, and explain why exceptions should become rare.
\item The number $L(x)$ is positive on a stretch near $9.06\times10^8$. Find the next stretches (they are known from the literature on P\'olya's problem) and
compare them with the predictions of the explicit formula.
\item Does Euler's sum of powers conjecture hold for $k=6$? No counterexample is known: no six sixth powers, or fewer, are known to sum to a sixth power.
\item Design your own ``Borwein'' sequence, with $a_k$ whose partial sums cross $1$ as late as you like, and compute the deficit exactly.
\end{itemize}

\section{Exercises}

\begin{exercise}[$\star$]\label{conjectures:ex:euler41}
Check that $n^2+n+41$ is prime for $0\le n\le39$ and not for $n=40$. Show that no non-constant polynomial $f$ with integer coefficients is prime for every
integer $n\ge0$.
\end{exercise}

\begin{exercise}[$\star$]\label{conjectures:ex:fermat641}
Using $641=5\cdot2^7+1=5^4+2^4$, show that $641$ divides $2^{32}+1$, without multiplying large numbers.
\end{exercise}

\begin{exercise}[$\star$]\label{conjectures:ex:threes}
Let $N_k=33\ldots31$, with $k$ threes. Show that $N_k=\frac{10^{k+1}-7}3$, check that $17\mid N_8$, and deduce that $N_k$ is composite for infinitely many $k$.
\end{exercise}

\begin{exercise}[$\star\star$]\label{conjectures:ex:moser}
Put $n$ points on a circle so that no three chords meet at one interior point, and draw all chords. Show that the number of regions inside the circle is
$\binom n4+\binom n2+1$, and compute it for $n=1,\dots,7$.
\end{exercise}

\begin{exercise}[$\star$]\label{conjectures:ex:squares}
Show that $\sum_{d\mid n}\lambda(d)$ is $1$ if $n$ is a square and $0$ otherwise. Deduce that $\sum_{n\le x}\lambda(n)\bigl\lfloor\frac xn\bigr\rfloor=\lfloor\sqrt x\rfloor$.
\end{exercise}

\begin{exercise}[$\star$]\label{conjectures:ex:squarefree}
Show that $\sum_{n\ge1}\mu(n)^2n^{-s}=\frac{\zeta(s)}{\zeta(2s)}$ and $\sum_{n\ge1}2^{\omega(n)}n^{-s}=\frac{\zeta(s)^2}{\zeta(2s)}$ for $\operatorname{Re}s>1$, where
$\omega(n)$ is the number of distinct prime factors of $n$.
\end{exercise}

\begin{exercise}[$\star\star$]\label{conjectures:ex:etaentire}
Show that $\eta(s)=\sum_{n\ge1}(-1)^{n-1}n^{-s}$ converges for real $s>0$ and diverges for real $s\le0$, and that $\eta(s)=(1-2^{1-s})\zeta(s)$ extends to an entire
function. Why does this not contradict Theorem~\ref{conjectures:thm:landau}?
\end{exercise}

\begin{exercise}[$\star\star$]\label{conjectures:ex:Lnegative}
Show, without any hypothesis, that $L(x)<0$ for arbitrarily large $x$.
\end{exercise}

\begin{exercise}[$\star\star$]\label{conjectures:ex:mertenssign}
Show that $M(x)=\sum_{n\le x}\mu(n)$ changes sign infinitely often. You may use that $\zeta$ has a zero at $\frac12+14.1347\ldots i$.
\end{exercise}

\begin{exercise}[$\star$]\label{conjectures:ex:residue}
Show that $\frac{\zeta(2s)}{s\zeta(s)}$ has a simple pole at $s=\frac12$ with residue $\frac1{\zeta(1/2)}$.
\end{exercise}

\begin{exercise}[$\star$]\label{conjectures:ex:sincpairs}
Use Theorem~\ref{conjectures:thm:borwein} to compute $\int_0^\infty\operatorname{sinc}^2x\dd x$, $\int_0^\infty\operatorname{sinc}x\operatorname{sinc}\frac x2\dd x$ and
$\int_0^\infty\operatorname{sinc}x\operatorname{sinc}2x\dd x$.
\end{exercise}

\begin{exercise}[$\star\star$]\label{conjectures:ex:harmonicsinc}
Show that $\int_0^\infty\operatorname{sinc}x\operatorname{sinc}\frac x2\cdots\operatorname{sinc}\frac x{n}\dd x=\frac\pi2$ for $n\le3$, and that for $n=4$ it equals
$\frac\pi2-\frac\pi{3456}$.
\end{exercise}

\section{Hints}

\begin{hint}{conjectures:ex:euler41}
For the second part, look at $f(kp)$ where $p=f(0)$, or $f(m)$ for a well-chosen $m$ if $f(0)\in\{0,\pm1\}$.
\end{hint}
\begin{hint}{conjectures:ex:fermat641}
Work modulo $641$: $5\cdot2^7\equiv-1$, so $5^4\cdot2^{28}\equiv1$.
\end{hint}
\begin{hint}{conjectures:ex:threes}
The order of $10$ modulo $17$ is $16$.
\end{hint}
\begin{hint}{conjectures:ex:moser}
Count the vertices and edges of the planar graph formed by the circle, the points, the chords and their intersections, and use Euler's formula.
\end{hint}
\begin{hint}{conjectures:ex:squares}
The function $\sum_{d\mid n}\lambda(d)$ is multiplicative: compute it on prime powers.
\end{hint}
\begin{hint}{conjectures:ex:squarefree}
Euler products.
\end{hint}
\begin{hint}{conjectures:ex:etaentire}
The functional equation of Volume~5 continues $\zeta$ to $\C\setminus\{1\}$. What are the coefficients of $\eta$?
\end{hint}
\begin{hint}{conjectures:ex:Lnegative}
If $L(x)\ge0$ for large $x$, look at $\int_{x_0}^\infty L(x)x^{-s-1}\dd x$ as $s\to\frac12^+$.
\end{hint}
\begin{hint}{conjectures:ex:mertenssign}
Follow the proof of Theorem~\ref{conjectures:thm:polyarh} with $\frac1{s\zeta(s)}$, and look at the real point $0$.
\end{hint}
\begin{hint}{conjectures:ex:residue}
$\zeta(2s)=\frac1{2s-1}+O(1)$ near $s=\frac12$.
\end{hint}
\begin{hint}{conjectures:ex:sincpairs}
The sum $a_1$ is $1$, $\frac12$ and $2$.
\end{hint}
\begin{hint}{conjectures:ex:harmonicsinc}
Copy the proof of Corollary~\ref{conjectures:cor:borweindeficit}.
\end{hint}

\section{Solutions}

\begin{solution}{conjectures:ex:euler41}
\idea A polynomial value $f(m)$ divides $f(m+kf(m))-f(m)$, so a prime value $p$ recurs as a factor.
\writeup The first part is a check. For $n=40$, $40^2+40+41=40\cdot41+41=41^2$. Now let $f$ be non-constant with integer coefficients.
Since $f(m)-f(n)$ is divisible by $m-n$, for any $m$ and $k$ the number $f(m+kf(m))$ is divisible by $f(m)$. As $f$ is non-constant, $|f(n)|\to\infty$, so choose
$m\ge0$ with $|f(m)|\ge2$. If $f(m)$ is not prime we are done. If $f(m)=\pm p$ with $p$ prime, then $p$ divides $f(m+kp)$ for every $k\ge0$, and for $k$ large
$|f(m+kp)|>p$, so $f(m+kp)$ is not prime.
\end{solution}

\begin{solution}{conjectures:ex:fermat641}
\idea Make $2^{32}$ appear from the two expressions of $641$.
\writeup Modulo $641$, $5\cdot2^7\equiv-1$, so $5^4\cdot2^{28}\equiv(-1)^4=1$. Also $5^4\equiv-2^4$. Hence $-2^4\cdot2^{28}\equiv1$, that is $2^{32}\equiv-1$, and
$641\mid2^{32}+1$.
\end{solution}

\begin{solution}{conjectures:ex:threes}
\idea Numbers with a repeated digit are geometric sums, and divisibility by $17$ is periodic in $k$.
\writeup $N_k=3(10^k+\dots+10)+1=\frac{10^{k+1}-10}3+1=\frac{10^{k+1}-7}3$. For $k=8$, $N_8=333\,333\,331=17\times19\,607\,843$. Since $3$ is invertible mod $17$, $17\mid N_k$
if and only if $10^{k+1}\equiv7\pmod{17}$, and $10^{16}\equiv1\pmod{17}$ (Fermat), so $17\mid N_k$ for every $k\equiv8\pmod{16}$. These $N_k$ exceed $17$, so they are
composite.
\end{solution}

\begin{solution}{conjectures:ex:moser}
\idea Every interior crossing comes from a choice of $4$ points, and Euler's formula turns vertices and edges into faces.
\writeup Each set of $4$ points gives exactly one crossing of two chords (the diagonals of their quadrilateral), so there are $I=\binom n4$ interior crossings.
Consider the planar graph whose vertices are the $n$ points and the $I$ crossings. A chord crossed by $c$ others is cut into $c+1$ edges; each crossing lies
on two chords, so the chords give $\binom n2+2I$ edges, and the circle gives $n$ arcs. Euler's formula $V-E+F=2$ gives
$F=2-n-I+\binom n2+2I+n=2+\binom n2+I$, and removing the outer face leaves $\binom n4+\binom n2+1$ regions. For $n=1,\dots,7$: $1,2,4,8,16,31,57$.
\end{solution}

\begin{solution}{conjectures:ex:squares}
\idea A multiplicative function is determined by prime powers.
\writeup $g(n)=\sum_{d\mid n}\lambda(d)$ is multiplicative as a convolution of multiplicative functions, and $g(p^k)=1-1+1-\cdots\pm1$ ($k+1$ terms) is $1$ if $k$ is
even and $0$ if $k$ is odd. So $g(n)=1$ exactly when all exponents are even. Then $\sum_{m\le x}g(m)=\sum_{m\le x}\sum_{d\mid m}\lambda(d)=\sum_{d\le x}\lambda(d)\lfloor\frac xd\rfloor$,
and the left side counts the squares up to $x$, which is $\lfloor\sqrt x\rfloor$.
\end{solution}

\begin{solution}{conjectures:ex:squarefree}
\idea Read off the local factor at each prime.
\writeup The local factor of $\mu^2$ is $1+p^{-s}=\frac{1-p^{-2s}}{1-p^{-s}}$, whose product is $\frac{\zeta(s)}{\zeta(2s)}$. The local factor of $2^\omega$ is
$1+2\sum_{k\ge1}p^{-ks}=\frac{1+p^{-s}}{1-p^{-s}}=\frac{1-p^{-2s}}{(1-p^{-s})^2}$, whose product is $\frac{\zeta(s)^2}{\zeta(2s)}$. The products converge absolutely for
$\operatorname{Re}s>1$ (Volume~5, Theorem~12.3).
\end{solution}

\begin{solution}{conjectures:ex:etaentire}
\idea Cancellation can make a series converge without making its function singular where convergence stops.
\writeup For real $s>0$ the terms alternate and decrease to $0$ in absolute value, so the series converges; for $s\le0$ the terms do not tend to $0$. On
$\operatorname{Re}s>1$, $\eta=(1-2^{1-s})\zeta$, and $\zeta$ is meromorphic on $\C$ with a single simple pole at $1$ (Volume~5, Theorem~12.8). The factor $1-2^{1-s}$
vanishes at $s=1$, so $(1-2^{1-s})\zeta(s)$ is entire. There is no contradiction, because the coefficients $(-1)^{n-1}$ do not have constant sign. Landau's theorem
does apply to the integral form: the summatory function $A(x)=\sum_{n\le x}(-1)^{n-1}$ takes only the values $0$ and $1$, so it is nonnegative, and
$\int_1^\infty A(x)x^{-s-1}\dd x=\frac{\eta(s)}s$ converges exactly for $\operatorname{Re}s>0$. As the theorem predicts, this function is singular at $0$: it has a
pole there, since $\eta(0)=\frac12\ne0$. Multiplying by $s$ removes the pole, which is why $\eta$ itself is entire.
\end{solution}

\begin{solution}{conjectures:ex:Lnegative}
\idea A positive integrand cannot produce an integral that tends to $-\infty$.
\writeup Suppose $L(x)\ge0$ for $x\ge x_0$. For real $s>1$,
\[
\frac{\zeta(2s)}{s\zeta(s)}=\int_1^{x_0}L(x)x^{-s-1}\dd x+\int_{x_0}^\infty L(x)x^{-s-1}\dd x\ge-C,
\]
with $C=\int_1^{x_0}|L(x)|\dd x$. By Theorem~\ref{conjectures:thm:landau} applied to $A=L$, exactly as in the proof of Theorem~\ref{conjectures:thm:polyarh}, the integral converges for
$s>\frac12$ and equals $\frac{\zeta(2s)}{s\zeta(s)}$ there, so the inequality holds for real $s>\frac12$. But as $s\to\frac12^+$, $\zeta(2s)\to+\infty$ while
$s\zeta(s)\to\frac12\zeta(\frac12)<0$ (Lemma~\ref{conjectures:lem:zetaneg}), so the left side tends to $-\infty$. Contradiction.
\end{solution}

\begin{solution}{conjectures:ex:mertenssign}
\idea The function $\frac1{s\zeta(s)}$ is regular on the whole real half-line $]0,\infty[$, so Landau would push convergence down to $0$, past the zeros on the
critical line.
\writeup As for $L$, $\frac1{\zeta(s)}=\sum\mu(n)n^{-s}=s\int_1^\infty M(x)x^{-s-1}\dd x$ for $\operatorname{Re}s>1$. Suppose $M(x)\ge0$ for $x\ge x_0$ (the case
$M\le0$ is the same with $-M$). The function $\frac1{s\zeta(s)}$ is holomorphic near every real point $s>0$: $\zeta$ has no zero on $]0,\infty[$ (Lemma~\ref{conjectures:lem:zetaneg}
and the Euler product), and $\frac1\zeta$ vanishes at the pole $s=1$. By Theorem~\ref{conjectures:thm:landau}, as in the proof of Theorem~\ref{conjectures:thm:polyarh}, the integral
converges and is holomorphic on $\operatorname{Re}s>0$, and equals $\frac1{s\zeta(s)}$ there away from the zeros of $\zeta$. But $\frac1{s\zeta(s)}$ has a pole at
$\frac12+14.1347\ldots i$, near which the integral would be bounded. Contradiction.
\end{solution}

\begin{solution}{conjectures:ex:residue}
\idea Only $\zeta(2s)$ is singular at $\frac12$.
\writeup By Volume~5, Proposition~12.6, $\zeta(w)=\frac1{w-1}+O(1)$ near $w=1$, so $\zeta(2s)=\frac1{2(s-\frac12)}+O(1)$ near $s=\frac12$. The factor $\frac1{s\zeta(s)}$
is holomorphic at $\frac12$ with value $\frac2{\zeta(1/2)}$, since $\zeta(\frac12)\ne0$ (Lemma~\ref{conjectures:lem:zetaneg}). The residue is $\frac12\cdot\frac2{\zeta(1/2)}=\frac1{\zeta(1/2)}=-0.6848\ldots$
\end{solution}

\begin{solution}{conjectures:ex:sincpairs}
\idea Each integral is $\frac\pi2\mathbb P(|a_1U|<1)$.
\writeup With $a_1=1$, $|U|<1$ almost surely: $\int_0^\infty\operatorname{sinc}^2x\dd x=\frac\pi2$. With $a_1=\frac12$, also $\frac\pi2$. With $a_1=2$,
$\mathbb P(|2U|<1)=\mathbb P(|U|<\frac12)=\frac12$, so $\int_0^\infty\operatorname{sinc}x\operatorname{sinc}2x\dd x=\frac\pi4$.
\end{solution}

\begin{solution}{conjectures:ex:harmonicsinc}
\idea The partial sums $\frac12$, $\frac56$ of $\frac12+\frac13+\cdots$ stay below $1$, and $\frac12+\frac13+\frac14=\frac{13}{12}$ crosses it by $\frac1{12}$.
\writeup Here $a_k=\frac1{k+1}$ for $k=1,\dots,n-1$. For $n\le3$, $\sum a_k\le\frac56<1$ and Theorem~\ref{conjectures:thm:borwein} gives $\frac\pi2$. For $n=4$,
$\varepsilon=\frac{13}{12}-1=\frac1{12}$, and the deficit is $\pi\,\mathbb P(S\ge1)$. With $V_k=1-U_k$ uniform on $[0,2]$, $S\ge1$ means $\sum a_kV_k\le\varepsilon$,
which forces $V_k\le\frac\varepsilon{a_k}\le\frac4{12}<2$; the simplex has volume $\frac{\varepsilon^3}{3!\,a_1a_2a_3}=\frac{24}{6\cdot1728}=\frac1{432}$ and the density is
$2^{-3}$, so $\mathbb P(S\ge1)=\frac1{3456}$ and the integral is $\frac\pi2-\frac\pi{3456}$. A numerical integration with mpmath confirms this to $15$ digits.
\end{solution}

\section*{Sources and further reading}

\begin{itemize}
\item R.~K.~Guy, \emph{The strong law of small numbers}, American Mathematical Monthly 95 (1988), 697--712 (the epigraph, Moser's circles, the threes and many
other examples); \emph{The second strong law of small numbers}, Mathematics Magazine 63 (1990), 3--20.
\item G.~P\'olya, \emph{Verschiedene Bemerkungen zur Zahlentheorie}, Jahresbericht der Deutschen Mathematiker-Vereinigung 28 (1919), 31--40.
C.~B.~Haselgrove, \emph{A disproof of a conjecture of P\'olya}, Mathematika 5 (1958), 141--145. R.~S.~Lehman, \emph{On Liouville's function}, Mathematics of
Computation 14 (1960), 311--320. M.~Tanaka, \emph{A numerical investigation on cumulative sum of the Liouville function}, Tokyo Journal of Mathematics 3 (1980),
187--189.
\item E.~Landau, \emph{\"Uber einen Satz von Tschebyschef}, Mathematische Annalen 61 (1905), 527--550; H.~L.~Montgomery and R.~C.~Vaughan, \emph{Multiplicative Number Theory I}, Cambridge University Press, 2007, Section~1.3 (Landau's theorem).
\item A.~M.~Odlyzko and H.~J.~J.~te Riele, \emph{Disproof of the Mertens conjecture}, Journal f\"ur die reine und angewandte Mathematik 357 (1985), 138--160
(the Stieltjes letter and the bounds $1.06$ and $-1.009$). J.~Pintz, \emph{An effective disproof of the Mertens conjecture}, Ast\'erisque 147--148 (1987),
325--333. G.~Hurst, \emph{Computations of the Mertens function and improved bounds on the Mertens conjecture}, Mathematics of Computation 87 (2018), arXiv:1610.08551.
\item C.~Bays and R.~H.~Hudson, \emph{A new bound for the smallest $x$ with $\pi(x)>\operatorname{li}(x)$}, Mathematics of Computation 69 (2000), 1285--1296.
\item L.~Hodges, \emph{A lesser-known Goldbach conjecture}, Mathematics Magazine 66 (1993), 45--47 (Goldbach's letter of 18 November 1752 and Stern's 1856
counterexamples, published in M.~A.~Stern, Nouvelles Annales de Math\'ematiques 15 (1856), 23--24).
\item L.~J.~Lander and T.~R.~Parkin, \emph{Counterexample to Euler's conjecture on sums of like powers}, Bulletin of the American Mathematical Society 72 (1966),
1079. N.~D.~Elkies, \emph{On $A^4+B^4+C^4=D^4$}, Mathematics of Computation 51 (1988), 825--835. R.~E.~Frye, \emph{Finding $95800^4+217519^4+414560^4=422481^4$ on the Connection Machine}, Proceedings of Supercomputing~88 (1988).
\item M.~K\v{r}\'{\i}\v{z}ek, F.~Luca and L.~Somer, \emph{17 Lectures on Fermat Numbers}, Springer, 2001 (Fermat's belief and Euler's factor $641$).
\item D.~Borwein and J.~M.~Borwein, \emph{Some remarkable properties of sinc and related integrals}, The Ramanujan Journal 5 (2001), 73--89 (the exact value of
the eighth integral); the probabilistic proof follows the well-known interpretation of $\operatorname{sinc}$ as the characteristic function of a uniform variable.
\end{itemize}
